\textbf{Neutronics}

From a neutronics perspective, Last Breath Reactor has proven the capability of generating thermal energy from spent nuclear fuel rods while simultaneously reducing the quality of plutonium. With the use of the \textit{Power Inverter}, LBR can sustain criticality up to an averaged burnup of 61.6 GWd. However, the Power Inverter introduces a proliferation risk -- containg weapons grade plutonium at the end of cycle -- and diverts from the central goal of avoiding fresh fuel usage. A better approach might be optimizing a fuel reloading pattern such that burnup can be even further extended without relying on fresh fuel. This is a complex optimization problem that must consider fission project generation, structural integrity, and neutron economy.


\textbf{Structural \& Operational}
The 6-cm lattice pitch will make the reactor considerably large. To mitigate the sheer size, the reactor should operate with a power density that is large, but reasonable enough to meet materials and thermal-hydraulics safety margins. With a higher power density comes faster depletion -- to make operation economical, the reactor vessel should be designed for on-line reloading. Some considerations should be made about vessel orientation, the use of pressure tubes (like a CANDU), etc.

